% Databases & Query Optimisation section

% 1 – SQL Server vs Oracle vs PostgreSQL
\QuestionSlide[\CategoryBadge[DbColor!20]{Databases}]{* What are the key differences between SQL Server, Oracle, and PostgreSQL?}

\begin{frame}[fragile]
  \frametitle{%
    \begin{tikzpicture}[remember picture, overlay]
      \node[anchor=north east, xshift=-0.4cm, yshift=-0.4cm, text=black] at (current page.north east) {
        \CategoryBadge[DbColor!20]{Databases}
      };
    \end{tikzpicture}
    Answer \theqcounter: SQL Server vs Oracle vs PostgreSQL?%
  }

  {\footnotesize
  \begin{tabular}{p{2.8cm}p{2.8cm}p{2.8cm}}
    \textbf{SQL Server} & \textbf{Oracle DB} & \textbf{PostgreSQL} \\
    T-SQL dialect & PL/SQL dialect & PL/pgSQL (SQL standard) \\
    Windows-first, Linux since 2017 & Cross-platform & Cross-platform, open-source \\
    Deep Azure/AD integration & Enterprise RAC clustering & JSONB, extensions, foreign data wrappers \\
    Azure SQL, RDS & Exadata, OCI & Neon, Supabase, Aurora \\
  \end{tabular}
  }
\end{frame}

% 2 – Indexes
\QuestionSlide[\CategoryBadge[DbColor!20]{Databases}]{** What are database indexes and their key types?}

\begin{frame}[fragile]
  \frametitle{%
    \begin{tikzpicture}[remember picture, overlay]
      \node[anchor=north east, xshift=-0.4cm, yshift=-0.4cm, text=black] at (current page.north east) {
        \CategoryBadge[DbColor!20]{Databases}
      };
    \end{tikzpicture}
    Answer \theqcounter: What are indexes and their types?%
  }

  {\footnotesize
  An \textbf{index} is a data structure that speeds up row lookups at the cost of write overhead and storage. Key types:
  \begin{itemize}
    \item \textbf{Clustered} (SQL Server) – defines physical row order; one per table.
    \item \textbf{Non-clustered} – separate B-tree with row pointers; multiple allowed.
    \item \textbf{Covering index} – includes all queried columns to avoid key lookup.
    \item \textbf{Filtered index} – built over a WHERE predicate, smaller and faster.
    \item \textbf{Columnstore} – columnar storage for analytics / data warehouse queries.
  \end{itemize}
  }

  \begin{minted}[fontsize=\scriptsize]{sql}
-- Covering index: no key-lookup needed for this query
CREATE INDEX IX_Orders_CustomerId_Status
    ON Orders (CustomerId)
    INCLUDE (Status, TotalAmount);
  \end{minted}
\end{frame}

% 3 – Execution Plan
\QuestionSlide[\CategoryBadge[DbColor!20]{Databases}]{** What is an execution plan and how do you use it?}

\begin{frame}[fragile]
  \frametitle{%
    \begin{tikzpicture}[remember picture, overlay]
      \node[anchor=north east, xshift=-0.4cm, yshift=-0.4cm, text=black] at (current page.north east) {
        \CategoryBadge[DbColor!20]{Databases}
      };
    \end{tikzpicture}
    Answer \theqcounter: What is an execution plan?%
  }

  {\footnotesize
  An \textbf{execution plan} is the query optimizer's step-by-step strategy for retrieving data. In SQL Server: \texttt{SET STATISTICS IO ON} or the graphical plan in SSMS. Look for: \textbf{Table Scan} (no index — bad), \textbf{Key Lookup} (missing include columns), \textbf{Sort} (no supporting index), and thick arrows (high row estimates). Use \texttt{EXPLAIN ANALYZE} in PostgreSQL.
  }

  \begin{minted}[fontsize=\scriptsize]{sql}
-- SQL Server: show actual execution plan
SET STATISTICS IO, TIME ON;

SELECT o.Id, o.Total
FROM   Orders o
WHERE  o.CustomerId = 42
  AND  o.Status = 'Pending';

-- Check: Index Seek (good) vs Table Scan (add index)
  \end{minted}
\end{frame}

% 4 – CTEs
\QuestionSlide[\CategoryBadge[DbColor!20]{Databases}]{** What is a CTE (Common Table Expression)?}

\begin{frame}[fragile]
  \frametitle{%
    \begin{tikzpicture}[remember picture, overlay]
      \node[anchor=north east, xshift=-0.4cm, yshift=-0.4cm, text=black] at (current page.north east) {
        \CategoryBadge[DbColor!20]{Databases}
      };
    \end{tikzpicture}
    Answer \theqcounter: What is a CTE?%
  }

  {\footnotesize
  A \textbf{CTE} (\texttt{WITH} clause) is a named temporary result set scoped to a single query. CTEs improve readability, allow self-referencing for \textbf{recursive queries} (hierarchies, graphs), and can be referenced multiple times in the same query. They are not materialised by default in SQL Server.
  }

  \begin{minted}[fontsize=\scriptsize]{sql}
-- Recursive CTE: employee hierarchy
WITH OrgTree AS (
    SELECT Id, Name, ManagerId, 0 AS Level
    FROM   Employees
    WHERE  ManagerId IS NULL         -- root

    UNION ALL

    SELECT e.Id, e.Name, e.ManagerId, t.Level + 1
    FROM   Employees e
    JOIN   OrgTree t ON e.ManagerId = t.Id
)
SELECT * FROM OrgTree ORDER BY Level;
  \end{minted}
\end{frame}

% 5 – Window Functions
\QuestionSlide[\CategoryBadge[DbColor!20]{Databases}]{** What are window functions in SQL?}

\begin{frame}[fragile]
  \frametitle{%
    \begin{tikzpicture}[remember picture, overlay]
      \node[anchor=north east, xshift=-0.4cm, yshift=-0.4cm, text=black] at (current page.north east) {
        \CategoryBadge[DbColor!20]{Databases}
      };
    \end{tikzpicture}
    Answer \theqcounter: What are window functions?%
  }

  {\footnotesize
  \textbf{Window functions} apply aggregate-like calculations over a \textbf{sliding window} of rows relative to the current row, without collapsing the result set. Key functions:
  \begin{itemize}
    \item \textbf{ROW\_NUMBER()} – unique sequential number per partition.
    \item \textbf{RANK() / DENSE\_RANK()} – rank with/without gaps on ties.
    \item \textbf{LAG() / LEAD()} – access previous/next row values.
    \item \textbf{SUM() OVER()} – running/cumulative totals.
  \end{itemize}
  }

  \begin{minted}[fontsize=\scriptsize]{sql}
SELECT CustomerId, OrderDate, Total,
       SUM(Total) OVER (PARTITION BY CustomerId
                        ORDER BY OrderDate
                        ROWS UNBOUNDED PRECEDING) AS RunningTotal,
       LAG(Total) OVER (PARTITION BY CustomerId
                        ORDER BY OrderDate) AS PrevOrderTotal
FROM Orders;
  \end{minted}
\end{frame}

% 6 – N+1 Problem
\QuestionSlide[\CategoryBadge[DbColor!20]{Databases}]{** What is the N+1 query problem?}

\begin{frame}[fragile]
  \frametitle{%
    \begin{tikzpicture}[remember picture, overlay]
      \node[anchor=north east, xshift=-0.4cm, yshift=-0.4cm, text=black] at (current page.north east) {
        \CategoryBadge[DbColor!20]{Databases}
      };
    \end{tikzpicture}
    Answer \theqcounter: What is the N+1 problem?%
  }

  {\footnotesize
  The \textbf{N+1 problem} occurs when loading a list of N entities triggers N additional queries to load each entity's related data — totalling N+1 round trips to the database. Fix in EF Core: use \texttt{.Include()} (eager loading) or a single \texttt{JOIN}. Detect with MiniProfiler or EF Core logging.
  }

  \begin{minted}[fontsize=\scriptsize]{csharp}
// BAD: 1 query for orders + N queries for each customer
var orders = await db.Orders.ToListAsync();
foreach (var o in orders)
    Console.WriteLine(o.Customer.Name); // lazy load per row

// GOOD: 1 query with JOIN
var orders = await db.Orders
    .Include(o => o.Customer)
    .ToListAsync();
  \end{minted}
\end{frame}

% 7 – Data Modelling Normal Forms
\QuestionSlide[\CategoryBadge[DbColor!20]{Databases}]{** What are the first three normal forms in data modelling?}

\begin{frame}[fragile]
  \frametitle{%
    \begin{tikzpicture}[remember picture, overlay]
      \node[anchor=north east, xshift=-0.4cm, yshift=-0.4cm, text=black] at (current page.north east) {
        \CategoryBadge[DbColor!20]{Databases}
      };
    \end{tikzpicture}
    Answer \theqcounter: What are the first three normal forms?%
  }

  {\footnotesize
  \begin{itemize}
    \item \textbf{1NF} – atomic values; no repeating groups or arrays in columns.
    \item \textbf{2NF} – 1NF + every non-key column depends on the \emph{whole} primary key (eliminates partial dependency; mainly relevant for composite keys).
    \item \textbf{3NF} – 2NF + no transitive dependencies (non-key columns depend only on the PK, not on other non-key columns).
  \end{itemize}
  \textbf{BCNF}, 4NF, 5NF address edge cases. In practice, 3NF is the standard target; read-heavy systems often deliberately \emph{denormalise} for performance.
  }
\end{frame}

% 8 – Query Optimisation Tips
\QuestionSlide[\CategoryBadge[DbColor!20]{Databases}]{*** What are the most impactful SQL query optimisation techniques?}

\begin{frame}[fragile]
  \frametitle{%
    \begin{tikzpicture}[remember picture, overlay]
      \node[anchor=north east, xshift=-0.4cm, yshift=-0.4cm, text=black] at (current page.north east) {
        \CategoryBadge[DbColor!20]{Databases}
      };
    \end{tikzpicture}
    Answer \theqcounter: Key query optimisation techniques?%
  }

  {\footnotesize
  \begin{itemize}
    \item \textbf{Index} the most selective columns; use covering indexes to avoid key lookups.
    \item \textbf{Avoid functions on indexed columns} in WHERE — they prevent index seek: \texttt{YEAR(OrderDate) = 2025} $\to$ \texttt{OrderDate >= '2025-01-01'}.
    \item \textbf{Use SET-based} operations instead of row-by-row cursors.
    \item \textbf{Paginate} with \texttt{OFFSET/FETCH} or keyset pagination for large result sets.
    \item \textbf{Update statistics} and rebuild fragmented indexes regularly.
    \item \textbf{Read execution plans} to spot scans, sorts, and spills.
  \end{itemize}
  }
\end{frame}
